\chapter{Thoracentesis}

\section*{Overview}
Thoracentesis is a procedure that removes fluid from the space between the lung and chest wall (pleural space) through a needle or catheter, for diagnosis and symptom relief.

\section*{Alternative Names}
Pleural tap; pleural fluid aspiration; chest tap

\section*{Principle}
A needle is advanced into the pleural space, usually under ultrasound guidance, to aspirate fluid. Fluid analysis distinguishes transudates from exudates and may identify infection, malignancy, or other causes.

\section*{History}
Thoracentesis has been performed for over a century, and ultrasound guidance, introduced in the late 20th century, greatly reduced complication rates.

\section*{Why It Is Done}
Performed for new or unexplained pleural effusion to determine its cause, and to relieve dyspnea from large effusions.

\section*{Is This a Primary or Secondary Test or Procedure}
Primary diagnostic and therapeutic procedure for pleural effusion.

\section*{How It Is Performed}
Performed at the bedside with the patient sitting upright. Ultrasound locates the fluid, the site is prepped and anesthetized, and a needle with or without a catheter is advanced above the rib over the fluid. Fluid is aspirated and sent for analysis.

\section*{Preparation}
Chest imaging (radiograph or ultrasound) to confirm fluid, coagulation assessment, and informed consent. Coughing or a ventilator may complicate the procedure.

\section*{Contraindications and Precautions}
Severe coagulopathy, thrombocytopenia, and overlying skin infection are relative contraindications. Precaution: avoid the diaphragm, and drain large volumes slowly to prevent re-expansion pulmonary edema.

\section*{Risks}
Pneumothorax, bleeding, infection, injury to the lung, liver, or spleen, and re-expansion pulmonary edema.

\section*{Aftercare and Recovery}
A chest radiograph is often performed to exclude pneumothorax, and the patient is observed briefly.

\section*{Results and Interpretation}
Pleural fluid protein, LDH, cell count, pH, cytology, and culture distinguish transudate (heart failure, cirrhosis) from exudate (infection, malignancy, inflammation).

\section*{Expected Results}
A small amount of clear fluid with low protein and cell counts, and relief of dyspnea.

\section*{Follow-up}
Results guide treatment; recurrent effusion may require repeat drainage, pleurodesis, or a chest tube.

\section*{References}
\begin{enumerate}
  \item MedlinePlus. Thoracentesis. U.S. National Library of Medicine; 2025.
  \item Current Medical Diagnosis and Treatment. McGraw-Hill; 2025.
\end{enumerate}

\clearpage
